\section{Introduction}
\label{sec:intro}

For two decades, the commercial web's surveillance economy operated on stateful storage primitives: HTTP cookies, Flash Local Shared Objects, and HTML5 Web Storage~\cite{eckersley2010panopticlick, englehardt2016openwpm}. In response to consumer pushback and regulatory mandates, major browser vendors introduced storage partitioning architectures, notably Apple WebKit's Intelligent Tracking Prevention (ITP) and Mozilla Gecko's Total Cookie Protection (dFPI). As third-party cookies decay, tracking consortiums have systematically pivoted to \textbf{stateless browser fingerprinting}: interrogating hardware, driver, and OS-level attributes through Web APIs (Canvas, WebGL, AudioContext, WebRTC, and system font enumeration) to reconstruct persistent cross-site behavioral graphs~\cite{laperdrix2016amiunique, iqbal2021fpinspector}.

To counteract stateless tracking, modern browsers deploy two fundamentally diverging anti-fingerprinting philosophies:
\begin{enumerate}
    \item \textbf{API Normalization:} Emulated by Tor Browser and Mozilla's Resist Fingerprinting (\texttt{privacy.resistFingerprinting}, RFP), which quantizes high-resolution timers, spoofs system attributes to standardized pools, and forces canvas readbacks to uniform white screens to maximize client anonymity sets~\cite{demir2022reproducibility}.
    \item \textbf{API Randomization (Farbling):} Implemented by Brave Browser, which injects cryptographically pseudo-random, session-ephemeral noise into canvas rendering, audio buffers, and hardware device enumeration, ensuring that subsequent readbacks generate divergent hashes that disrupt cross-session linking~\cite{liu2025fptrace}.
\end{enumerate}

Despite these deployments, the measurement literature lacks a rigorous methodology to determine whether these defenses functionally protect users in production. Current empirical benchmarks suffer from two structural flaws:

\paragraph{The Intent vs. Outcome Fallacy}
Passive measurement platforms (e.g., OpenWPM~\cite{englehardt2016openwpm}, FP-Inspector~\cite{iqbal2021fpinspector}) instrument JavaScript APIs to log when heuristic execution patterns occur. While effective for detecting tracking *intent*, observing a canvas readback provides zero verification of tracking *success*. It cannot determine whether the extracted payload resolved a user within an ad-tech identity graph or failed to overcome environmental noise.

\paragraph{The Theoretical Entropy Gap}
Synthetic benchmarks (e.g., Panopticlick~\cite{eckersley2010panopticlick}, AmIUnique~\cite{laperdrix2016amiunique}) benchmark defenses by calculating the Shannon entropy bits subtracted from isolated Web APIs. However, theoretical entropy reductions do not guarantee tracking severance in the wild. If an advertising network combines a partially farbled canvas signature with residual cross-site signals, such as IP-range localization, TLS client fingerprints, and User-Agent headers. Client re-identification persists despite nominal entropy loss.

\paragraph{Real-Time Bidding as an Economic Oracle}
To overcome the limitations of proxy measurements, we exploit a direct side-channel oracle: \textbf{programmatic Real-Time Bidding (RTB) clearing prices}~\cite{liu2025fptrace}. In header bidding auctions orchestrated via client-side Prebid.js wrappers, algorithmic Demand-Side Platforms (DSPs) submit bids (Cost Per Mille, or \cpm) within milliseconds based on identity certainty~\cite{pachilakis2019imc}. Econometric studies establish that an identified user commanding a rich cross-site behavioral profile sells for a $30\%$ to $60\%+$ valuation premium over an untargeted contextual impression~\cite{beales2014navigant, johnson2020marketsci, google2019cookies}. 

When an anti-fingerprinting defense successfully breaks cross-site tracking across an enforced stateful reset, DSP identity graphs are severed. Lacking identity certainty, bidding algorithms retreat to contextual reserve floors, precipitating a measurable collapse in \cpm clearing prices.

\paragraph{Contributions}
To quantify this phenomenon, we built a cluster-scale causal measurement platform that drives native Linux browser binaries across 150 commercial publishers. Our work delivers four primary contributions:
\begin{itemize}
    \item \textbf{Causal Measurement Methodology:} We design a 3-phase experimental protocol that trains high-intent personas, records pre-break ad-auction baselines, programmatically purges stateful storage, and measures post-purge economic bid shifts, causally isolating stateless tracking in the wild.
    \item \textbf{Validation of the Negative Control:} We demonstrate that in unhardened Google Chrome, stateful storage purges yield a negligible bid shift of only \textbf{$-1.12\%$} ($p = 0.4832$, fail to reject $H_0$). This provides empirical proof that clearing cookies without anti-fingerprinting defenses is entirely futile against commercial ad tech.
    \item \textbf{The Empirical Pareto Frontier:} We characterize an operational trade-off in browser engineering: \textit{Hardened Firefox} achieves the \textbf{Stealth / Zero-Breakage Optimum} ($-26.87\%$ median bid collapse, $p = 0.0734$, with zero usability breakage: $\ucost = 0.00$); whereas \textit{Hardened Brave} reaches the \textbf{Maximum Privacy Optimum} ($-59.97\%$ median collapse, $\rid = 0.500$), but imposes a $+44\%$ increase in runtime script exceptions ($\ucost = 0.44$).
    \item \textbf{Deconstructing the Usability Trap \& Metric Refinement:} We show that defensive request blocking triggers tracker retry loops ($28,976$ requests in Firefox Hardened) that invert naive evaluation metrics. We formalize the Refined Effectiveness-to-Risk metric ($\etr$), resolving structural pathologies and restoring monotonic rigor to privacy defense evaluation.
\end{itemize}